%! Special Directions --- Setting Up Eigenvalues — Unit 6, Lesson 6.0 — Lesson Plan
\documentclass[10pt]{article}
\usepackage{linalg-article}
\usepackage{linalg-boxes}
\usepackage{graphicx}   % -article does NOT load graphicx; needed for thumbnails/screenshots
\graphicspath{{images/}}
\newcommand{\TallMath}[1]{$\displaystyle #1\rule[-1.4em]{0pt}{3.2em}$}
\newcommand{\xx}{\mathbf{x}}
\newcommand{\yy}{\mathbf{y}}
\newcommand{\vv}{\mathbf{v}}
\newcommand{\zero}{\mathbf{0}}
\newcommand{\UnitNumberName}{Unit 6: Eigenvalues and Eigenvectors \quad}
\newcommand{\LessonNumberName}{Lesson 6.0: Special Directions --- Setting Up Eigenvalues}

\begin{document}
\begin{center}
    {\Large\bfseries \CourseName: \SchoolYear \\
    {\small\bfseries \UnitNumberName \LessonNumberName}}
\end{center}

\begin{tcolorbox}[colback=blush, colframe=burgundy, boxrule=0.9pt, arc=2mm,
    left=2mm, right=2mm, top=1.5mm, bottom=1.5mm]
    \textbf{Primary Objective:} Students can identify the \emph{special directions} of a matrix:
    a nonzero vector $\xx$ is an \emph{eigenvector} of $A$ when $A\xx$ lands on the same line as
    $\xx$ --- that is, $A\xx=\lambda\xx$ for a number $\lambda$ (the \emph{eigenvalue}), the factor
    by which that direction is stretched. Given a matrix and a candidate vector they can compute
    $A\xx$, decide whether it is a scalar multiple of $\xx$, and read off $\lambda$; and given an
    eigenvector they can find $\lambda$. They can explain \emph{why} these directions matter --- along
    them a matrix does nothing but \emph{scale}, so it acts as simply as a number --- and connect
    $\lambda=0$ to the Unit~5 fact $\det A=0$. They can describe where Unit~6 is going: \emph{finding}
    eigenvalues from $\det(A-\lambda I)=0$ (6.1), \emph{diagonalizing} $A$ (6.2), symmetric matrices
    (6.3), and an application (6.4).
\end{tcolorbox}

\renewcommand{\arraystretch}{1.4}
\begin{skillbox}[Priority Ideas \& Skills]{goldbox}
\begin{tabularx}{\linewidth}{@{}X|X@{}}
\textbf{Priority Ideas \& Skills} & \textbf{Key Understandings (the ``why'')} \\[2pt]
\hline
\vspace{-6pt}
\begin{itemize}[leftmargin=1.1em, nosep, after=\vspace{2pt}]
  \item Compute $A\xx$ and decide whether it is a \textbf{multiple} of $\xx$.
  \item State the eigen-relation $A\xx=\lambda\xx$; read $\lambda$ as the \textbf{stretch factor}.
  \item Given an eigenvector, \textbf{find} its eigenvalue $\lambda$.
  \item Connect $\lambda=0$ to $\det A=0$ (a squashed, non-invertible matrix --- Unit~5).
\end{itemize}
&
\vspace{-6pt}
\begin{itemize}[leftmargin=1.1em, nosep, after=\vspace{2pt}]
  \item A matrix usually \emph{turns} a vector; a few directions are only \emph{stretched}.
  \item Along an eigenvector a matrix behaves like a single number $\lambda$ --- no turning.
  \item The eigenvector is the \emph{direction}; the eigenvalue is \emph{how much} it scales.
  \item $\lambda=0$ means some direction is crushed to $\zero$ --- exactly the Unit~5 collapse.
\end{itemize}
\end{tabularx}
\end{skillbox}

\begin{skillbox}[{Vocabulary, Concepts, \& Theorems}]{greenbox}
\renewcommand{\arraystretch}{1.3}
\begin{tabularx}{\linewidth}{@{}l X@{}}
\textbf{Eigenvector} & A \emph{nonzero} vector $\xx$ whose image $A\xx$ lies on the same line as $\xx$ (direction unchanged, up to a flip). \\
\textbf{Eigenvalue $\lambda$} & The number with $A\xx=\lambda\xx$: the factor the eigenvector is stretched ($|\lambda|$) and flipped ($\lambda<0$) by. \\
\textbf{The eigen-relation} & $A\xx=\lambda\xx$ with $\xx\ne\zero$ --- ``matrix times vector $=$ number times the same vector.'' \\
\textbf{Test for an eigenvector} & Compute $A\xx$; if $A\xx=\lambda\xx$ for some $\lambda$, $\xx$ is an eigenvector and $\lambda$ is its eigenvalue. \\
\textbf{$\lambda=0$ \& singular} (recall §5) & $\lambda=0$ means $A\xx=\zero$ for a nonzero $\xx$ --- the matrix is \emph{not invertible} ($\det A=0$). \\
\textbf{Characteristic equation} (preview 6.1) & $\det(A-\lambda I)=0$ --- the equation whose solutions are the eigenvalues. \\
\end{tabularx}
\end{skillbox}

\begin{fixedskillbox}[Activate Prior Knowledge \& Spiral Review]{blush}
    The Warm-Up rehearses the three ideas this lesson builds on:
    \textbf{(1)} computing a \emph{matrix times a vector} $A\xx$ (Unit~1),
    \textbf{(2)} recognizing when one vector is a \emph{scalar multiple} of another (parallel /
    same line, Unit~1), and
    \textbf{(3)} the $2\times2$ \emph{determinant} $ad-bc$ from Unit~5 --- the tool that will later
    find the eigenvalues. Debrief items~1--2 together: $A\begin{bsmallmatrix}1\\1\end{bsmallmatrix}
    =\begin{bsmallmatrix}3\\3\end{bsmallmatrix}=3\begin{bsmallmatrix}1\\1\end{bsmallmatrix}$, so the
    matrix did nothing to that direction but \emph{multiply it by $3$}. That is today's whole idea.
\end{fixedskillbox}

\begin{skillbox}[Hook]{blush}
    Show a matrix acting on the plane (Lesson~5.3): most arrows get \emph{turned} to a new direction.
    Then ask: \textbf{``Are there any arrows the matrix does \emph{not} turn --- ones it only
    stretches?''} Test $\xx=\begin{bsmallmatrix}1\\1\end{bsmallmatrix}$ under
    $A=\begin{bsmallmatrix}2&1\\1&2\end{bsmallmatrix}$: $A\xx=\begin{bsmallmatrix}3\\3\end{bsmallmatrix}$
    --- same line, just three times as long. Name the payoff: for a few \emph{special directions} a
    matrix acts like a single \emph{number}. Pose the unit question: \textbf{``Which directions does a
    matrix only stretch, and by how much?''} Those directions are the \emph{eigenvectors}; the stretch
    factors are the \emph{eigenvalues}, and they run all of Unit~6 --- finding them (6.1), diagonalizing
    (6.2), symmetric matrices (6.3), and an application (6.4).
\end{skillbox}

\begin{skillbox}[Lesson]{blush}
\begin{multicols}{2}
\textbf{1. A matrix usually changes direction.} Recall $A\xx$ moves a vector (Lesson~5.3). For
$A=\begin{bsmallmatrix}2&1\\1&2\end{bsmallmatrix}$, the vector
$\begin{bsmallmatrix}1\\0\end{bsmallmatrix}$ maps to $\begin{bsmallmatrix}2\\1\end{bsmallmatrix}$ ---
a \emph{new} direction. Most vectors get turned.

\textbf{2. Special directions: $A\xx=\lambda\xx$.} A nonzero $\xx$ is an \textbf{eigenvector} when
$A\xx$ stays on the \emph{same line} --- $A\xx=\lambda\xx$. The number $\lambda$ is the
\textbf{eigenvalue}. Test $\begin{bsmallmatrix}1\\1\end{bsmallmatrix}$: $A\xx=
\begin{bsmallmatrix}3\\3\end{bsmallmatrix}=3\begin{bsmallmatrix}1\\1\end{bsmallmatrix}$, so
$\lambda=3$. Test $\begin{bsmallmatrix}1\\-1\end{bsmallmatrix}$: $A\xx=
\begin{bsmallmatrix}1\\-1\end{bsmallmatrix}=1\!\cdot\!\begin{bsmallmatrix}1\\-1\end{bsmallmatrix}$,
so $\lambda=1$.

\columnbreak

\textbf{3. How to test a candidate.} Compute $A\xx$; ask ``is it a \emph{multiple} of $\xx$?'' If yes,
that multiple is $\lambda$; if no, $\xx$ is not an eigenvector. Non-example: $A
\begin{bsmallmatrix}2\\1\end{bsmallmatrix}=\begin{bsmallmatrix}5\\4\end{bsmallmatrix}$ --- not a
multiple of $\begin{bsmallmatrix}2\\1\end{bsmallmatrix}$, so it is turned.

\textbf{4. Why it matters + road ahead.} Along an eigenvector a matrix does nothing but \emph{scale},
so it acts like the single number $\lambda$ --- that is what makes eigenvectors powerful. Special
case: $\lambda=0$ means $A\xx=\zero$ for a nonzero $\xx$, so $A$ is \emph{not invertible} ($\det A=0$,
Unit~5). \textbf{Finding} the $\lambda$'s comes next: they solve $\det(A-\lambda I)=0$ (a Unit~5
determinant $\Rightarrow$ a quadratic). \textbf{Road ahead:} \textbf{6.1} finding eigenvalues;
\textbf{6.2} diagonalizing; \textbf{6.3} symmetric matrices; \textbf{6.4} an application.
\end{multicols}
\end{skillbox}

\begin{skillbox}[Active Monitoring]{redbox}
Circulate during the notes practice and Tier work. Watch for and correct:
\begin{itemize}[leftmargin=1.2em, nosep]
  \item arithmetic slips in $A\xx$ (row $\cdot$ vector) --- the eigen-test is only as good as the product;
  \item declaring ``eigenvector'' without checking that $A\xx$ is \emph{the same multiple} in \emph{both} entries;
  \item forgetting the \emph{nonzero} requirement (the zero vector is never an eigenvector);
  \item thinking $\lambda$ must be positive --- $\lambda$ can be $0$ or negative (a flip).
\end{itemize}
Cold-call prompts: ``Compute $A\xx$ --- is it a multiple of $\xx$?'' ``What is $\lambda$ here?''
``If $A\xx=\zero$ for a nonzero $\xx$, what does that say about $\det A$?''
\end{skillbox}

\begin{skillbox}[Group Work \& Differentiation]{redbox}
\textbf{Finding the Special Directions} activity (tiered; mirrors the notes).
\begin{multicols}{3}
\textbf{Tier R --- Remediate.} Compute $A\xx$ for diagonal and given matrices; decide whether $A\xx$
is a multiple of $\xx$; read off $\lambda$ when it is.

\columnbreak
\textbf{Tier A --- Approaching.} Given eigenvectors, find each $\lambda$; identify which of several
vectors is the eigenvector; and meet a $\lambda=0$ eigenvector, connecting it to $\det A=0$ and
non-invertibility (recall §5).

\columnbreak
\textbf{Tier E --- Extension.} \emph{First look at finding eigenvalues:} form $A-\lambda I$, compute
the $2\times2$ determinant $\det(A-\lambda I)$, set it to $0$, and solve the quadratic --- recovering
the $\lambda$'s from the notes. This is the Lesson~6.1 method.
\end{multicols}
\end{skillbox}

\begin{skillbox}[Individual Work \& Assessment]{redbox}
\textbf{Exit Ticket} (independent, no notes): compute $A\xx$ and decide whether $\xx$ is an
eigenvector, giving $\lambda$; test a second vector that is \emph{not} an eigenvector; and a
\textbf{conceptual/justification check} --- state in one sentence what $A\xx=\lambda\xx$ means
geometrically. Collect and sort into ``got it / arithmetic slip / concept slip'' to decide whether
6.1 opens with a quick $A\xx$ re-warm.
\end{skillbox}

\begin{skillbox}[Reinforcement \& Extension]{goldbox}
\textbf{Homework:} compute $A\xx$ and test candidates; find $\lambda$ from given eigenvectors
(diagonal, triangular, symmetric); a $\lambda=0$ item tied to $\det A=0$ and invertibility (§5); and a
short justification of ``special direction $=$ only stretched.'' \textbf{Extension:} form $A-\lambda I$
and solve $\det(A-\lambda I)=0$ --- a first taste of Lesson~6.1. \textbf{Preview:} Lesson~6.1,
\emph{Introduction to Eigenvalues} --- how to \emph{find} the special directions (and their $\lambda$'s)
without guessing, from $\det(A-\lambda I)=0$.
\end{skillbox}
\end{document}
